% Autorregresión en el espacio tangente del baricentro de Wasserstein.
\begin{tikzpicture}[x=1cm,y=1cm]
% Superficie curva: espacio de distribuciones con la métrica W2
\fill[qazul!6] (0,0) .. controls (2.5,1.6) and (6.5,1.6) .. (9,0.2)
  .. controls (9.6,-1.2) and (9.2,-2.6) .. (8.4,-3.4)
  .. controls (6,-2.1) and (2.8,-2.2) .. (0.4,-3.2)
  .. controls (-0.3,-2.2) and (-0.4,-1) .. (0,0);
\draw[qazul!60,line width=0.6pt] (0,0) .. controls (2.5,1.6) and (6.5,1.6) .. (9,0.2)
  .. controls (9.6,-1.2) and (9.2,-2.6) .. (8.4,-3.4)
  .. controls (6,-2.1) and (2.8,-2.2) .. (0.4,-3.2)
  .. controls (-0.3,-2.2) and (-0.4,-1) .. (0,0);
\node[nota,anchor=west] at (0.3,-2.75) {distribuciones con la métrica $W_2$ (no lineal)};
% Plano tangente en el baricentro
\coordinate (B) at (4.6,-0.9);
\fill[qnaranja!10,draw=qnaranja!70,line width=0.6pt] (1.9,0.35) -- (7.9,0.35) -- (7.2,-2.05) -- (1.2,-2.05) -- cycle;
\node[nota,anchor=east,text=tinta] at (7.05,-1.8) {espacio tangente $T_{\bar{Q}}$ (lineal)};
\fill[tinta] (B) circle (2.2pt);
\node[font=\sffamily\scriptsize,anchor=north] at (B) {$\bar{Q}$};
% Log de dos sesiones y predicción
\draw[flecha,draw=qazuloscuro] (B) -- (2.6,-0.3) node[anchor=east,font=\sffamily\scriptsize]{$v_{t-1}$};
\draw[flecha,draw=qazuloscuro] (B) -- (3.4,0.05) node[anchor=south,font=\sffamily\scriptsize]{$v_t$};
\draw[flecha,draw=qnaranja,line width=0.9pt] (B) -- (5.4,0.1) node[anchor=south west,font=\sffamily\scriptsize]{$\hat v_{t+1}=A\,v_t$};
% Distribuciones en la superficie
\fill[qazul] (1.7,1.05) circle (2pt) node[anchor=south,font=\sffamily\scriptsize]{$Q_{t-1}$};
\fill[qazul] (3.2,1.25) circle (2pt) node[anchor=south,font=\sffamily\scriptsize]{$Q_t$};
\fill[qnaranja] (6.3,1.2) circle (2pt) node[anchor=south,font=\sffamily\scriptsize]{$\widehat{Q}_{t+1}$};
\draw[flecha,draw=tenue,densely dotted] (1.7,1.0) -- (2.55,-0.2);
\draw[flecha,draw=tenue,densely dotted] (3.2,1.2) -- (3.4,0.15);
\draw[flecha,draw=qnaranja,densely dotted] (5.45,0.2) -- (6.25,1.1);
\node[nota,anchor=west,align=left] at (9.4,1.0) {$\exp$ y proyección a\\ funciones cuantil crecientes};
\node[nota,anchor=east,align=right] at (-0.45,0.6) {$\log$ en una dimensión:\\ $v_t(u)=x_t(u)-\bar{x}(u)$};
\end{tikzpicture}
